\section{总体设计}

\subsection{体系结构}

系统采用 B/S 与消息机器人并行的三层结构（图~\ref{fig:arch}）。表示层是浏览器中的单页应用和 Telegram 客户端；应用层由三个进程组成——Web 服务、Telegram 机器人、后台任务——它们共享同一个业务服务层 \code{BillingService}，保证同一业务规则只有一份实现；数据层是计费数据库。用量库和网关属于外部系统：前者被只读挂载，后者通过 HTTP 管理接口访问。

\begin{figure}[H]
  \centering
  \includegraphics[width=\textwidth]{architecture.pdf}
  \caption{系统体系结构}\label{fig:arch}
\end{figure}

这种划分带来两点好处。第一，所有写操作都经过业务服务层并在同一事务中写入审计日志，无论入口是网页、机器人还是命令行。第二，数据库访问层只依赖 SQLAlchemy 的方言抽象，同一套代码既可运行在单机部署使用的 SQLite 上，也可运行在 PostgreSQL 上；数据库方言的差异集中在少数几处（见 6.2 节），由持续集成在两种数据库上分别跑完全部测试来保证一致。

\subsection{功能模块}

按照需求分析的 F1–F8，系统划分为七个功能模块（图~\ref{fig:modules}）。其中“用量采集计价”由后台任务驱动，其余模块由用户请求驱动。

\begin{figure}[H]
  \centering
  \includegraphics[width=\textwidth]{modules.pdf}
  \caption{功能模块结构}\label{fig:modules}
\end{figure}

\subsection{计费处理流程}

从一次调用到账单的一行，数据要经过图~\ref{fig:pipeline} 所示的处理链。前三步在调用发生后十几秒内由后台任务完成，结果写入计价记录；后五步在账期预览或关账时执行。

\begin{figure}[H]
  \centering
  \includegraphics[width=\textwidth]{pipeline.pdf}
  \caption{从调用记录到账单的处理流程}\label{fig:pipeline}
\end{figure}

\subsection{概念结构设计}

\subsubsection{实体识别}

从数据流图的数据存储出发，识别出如下实体：用户、API 密钥、调用记录、价格版本、模型价格、资源池、上游账号、上游分组、梯度规则、账期、账单、账单明细、账单调整。此外还有三个“带属性的联系”需要特别处理：

\begin{itemize}
  \item \textbf{归属期}：用户与密钥之间不是简单的一对多。一把密钥在不同时间段可以属于不同用户，联系本身带有“生效时间”“失效时间”两个属性。若只在密钥上记录“当前主人”，密钥转手后历史调用就无法正确归属。因此把它设计为联系，转换为关系时成为独立的表。
  \item \textbf{计价}：调用记录与价格版本之间是多对多联系。同一条调用在不同价格版本下各有一个计价结果，联系的属性是等效费用、计费档位和调用时刻的归属用户。
  \item \textbf{成本快照}：账期与上游账号之间是多对多联系。联系属性是该账号在该账期折算出的成本、汇率、分组名等，它们是关账时的历史事实。
\end{itemize}

\subsubsection{E-R 图}

图~\ref{fig:er} 给出系统的全局 E-R 图（陈氏表示法，矩形为实体，菱形为联系，椭圆为联系的属性；实体自身的属性较多，在第 4 章以关系模式列出）。

\begin{figure}[H]
  \centering
  \includegraphics[width=\textwidth]{er.pdf}
  \caption{全局 E-R 图}\label{fig:er}
\end{figure}

主要联系的语义与基数如下：
\begin{itemize}
  \item 用户—归属期—API 密钥（1:N，按时间）：同一时刻一把密钥至多属于一个用户，一个用户可以同时拥有多把密钥；
  \item API 密钥—使用—调用记录（1:N）：调用记录通过密钥摘要关联密钥，未登记的密钥产生的调用也要保留；
  \item 价格版本—包含—模型价格（1:N）：模型价格是弱实体，依赖价格版本存在；
  \item 调用记录—计价—价格版本（M:N）；计价结果归入一个资源池（N:1）；
  \item 账期—采用—价格版本（N:1）：账期创建时绑定生效版本，价格更新时开放账期跟随切换，已关账账期保持不变；
  \item 账期—出账—账单（1:N）且用户—出账—账单（1:N）：每个账期每位用户至多一张账单；
  \item 账单—组成—账单明细（1:N）；明细归入一个资源池（N:1）；
  \item 上游账号—归入—上游分组（N:1），上游分组—遵循—梯度规则（N:1）；
  \item 账期—调整—账单调整（1:N），调整针对某一用户。
\end{itemize}

\subsubsection{关键设计决策}

\paragraph{历史数据版本化而不覆盖}
价格、归属、成本、规则都会变化，而账单必须能复现，因此一律“追加新版本并引用版本号”：改价生成新价格版本，转手是关闭旧归属期、开启新归属期，关账时把规则与成本复制为快照。

\paragraph{计价与分摊分离}
计价只依赖单次调用和价格，可在调用后立即增量完成；分摊依赖整个账期，只在预览或关账时对计价记录聚合一次。

\paragraph{金额全部用整数}
货币量以 64 位整数存储（纳美元、分、百万分之一），既无浮点误差，也回避了 \code{NUMERIC} 在两种数据库上的差异；纳美元精度下可表示约 92 亿美元。
